\Needspace{24\baselineskip}\fancyhead[L]{\small COMET OF 1807.}
\centerline{COMET OF 1807.}

\begin{center}
\rule{0.5\linewidth}{0.4pt}\\[2pt]
Monatliche Correspondenz. Vol. XVI. Pp. 562–567. December 1807.\\[2pt]
\rule{0.5\linewidth}{0.4pt}
\end{center}
\par\indent ---\ ---\ --- Dr. \textsc{Gauss}, who pays attention to the orbit of every stranger that appears in the sky, wrote to us on 16 November. “From the Bremen and Lilienthal observations through 5 November, I derived the following parabolic elements of the comet

\smallskip

\begin{center}\begin{tabular}{lr}
Passage through perihelion, 18 Sept. .\ .\ .\ .\ . & $19^{\text{u}}\;6^{\text{m}}$ in Paris \\
Longitude of perihelion .\ .\ .\ .\ .\ .\ .\ .\ .\ .\ . & $271^\circ\;0'\;13''3$ \\
$\AscendingNode$ .\ .\ .\ .\ .\ .\ .\ .\ .\ .\ .\ .\ .\ .\ .\ . & $266\;38\;31.2$ \\
\multicolumn{2}{l}{both reckoned from the apparent equinox} \\
Log. of the least distance .\ .\ .\ .\ .\ . & $9.8114927$ \\
{}[Inclination of the orbit .\ .\ .\ .\ .\ .\ .\ .\ .\ . & $63^\circ\;12'\;35''6$] \\
\multicolumn{2}{l}{Prograde motion}
\end{tabular}\end{center}


\par\indent These elements differ very little from those I had previously calculated from a few Berlin observations. They likewise differ only slightly from those Mr. \textsc{Bessel} derived from a shorter series of observations. I will write to you more about this from Göttingen.”

\begin{flushright}
[\textsc{von Zach}.]
\end{flushright}

